\section{Setup, cost and full results}
\label{app:setup}
\begin{table}[h]
\centering\small
\resizebox{\linewidth}{!}{\input{../figures/tab_setup.tex}}
\caption{Cells: dimension $d$, instances explained, runs $R$ per instance, background size, test AUC and accuracy of the main model, reference type.}
\label{tab:setup}
\end{table}
\begin{table}[h]
\centering\small
\resizebox{\linewidth}{!}{\input{../figures/tab_cost.tex}}
\caption{Single-thread seconds per explainer call at the smallest and largest budget $B$ with the rows evaluated in parentheses. $B$ is the nominal budget (Kernel: coalitions; permutation: \texttt{max\_evals}), and the number of model rows actually evaluated differs between the explainers at the same $B$; timed on a shared, oversubscribed host, so indicative.}
\label{tab:cost}
\end{table}
\subsection{Smallest and largest budget per cell}
\begin{table}[h]
\centering\small
\resizebox{\linewidth}{!}{\input{../figures/c2_table_main.tex}}
\caption{Mean over instances of the probability $\bar p$ that one run's top-3 set equals the exact top-3 set, with 95\% instance-cluster bootstrap CI, and the fraction of instances for which this probability is $\ge0.9$, at the smallest and largest nominal budget $B$ of each cell (Kernel: coalitions; permutation: \texttt{max\_evals}); the Rows column gives the mean number of model rows evaluated per run, which differs between the explainers at the same $B$, so compare them at equal rows. Full grid: Appendix~\ref{app:full}.}
\label{tab:c2}
\end{table}
\subsection{Full top-$k$ match grid}\label{app:full}
{\scriptsize\setlength{\tabcolsep}{2.5pt}\input{../figures/c2_table_full.tex}}
\begin{figure}[h]
\centering\includegraphics[width=\linewidth]{fig_c4_diagnostic.pdf}
\caption{(Exploratory) AUC of per-instance diagnostics for flagging unreliable instances.}
\label{fig:c4}
\end{figure}
\subsection{Match probability against $z$, per cell}
\begin{figure}[h]
\centering\includegraphics[width=\linewidth]{fig_c2_z_cell.pdf}
\caption{Top-3 match probability against $z=$ reference gap$/$run-to-run sd of the gap, one line per model, 95\% CI, for each dataset and explainer. Dashed: the normal-approximation reference shape $\Phi(z)$ for one competing pair (computed by the plotting script, not fitted). The curves lie close to it because $z$ uses the sd of the same runs: the relation is largely built in and is not an independent test.}
\label{fig:zcell}
\end{figure}
\section{Default \texttt{l1\_reg}, reference check and log-odds}
\label{app:l1}
\begin{figure}[h]
\centering\includegraphics[width=\linewidth]{fig_l1_default.pdf}
\caption{German Credit (three models) and ACS Employment logistic regression, Kernel SHAP with \texttt{l1\_reg=False} (main grid) versus the library default \texttt{l1\_reg}: error and top-3 match against the reference (20\,000 coalitions for German Credit, exact for ACS Employment). The default variant plateaus because its target is the sparsified regression, a bias, not run-to-run noise.}
\label{fig:l1}
\end{figure}
\begin{table}[h]
\centering\scriptsize
\resizebox{\linewidth}{!}{\input{../figures/tab_retrain.tex}}
\caption{Retrain-subsample sensitivity: probability that two retrained models give the same top-3 set, for 50\%, 80\% (main protocol) and 95\% training subsamples (five retrains, 10 pairs per instance, 95\% CI over instances; the pairs are not independent). ``--'': not run (the 50\% and 95\% fractions were run for German Credit and ACS Income, and for ACS Employment only for logistic regression, forest and MLP). Exact explanations for $d\le16$, one Kernel run of 10\,000 coalitions for German Credit.}
\label{tab:retrain}
\end{table}
\begin{table}[h]
\centering\small
\input{../figures/tab_retrain_seed.tex}
\caption{German Credit: the estimator's contribution to retrain disagreement. Two 80\%-retrained models explained with one Kernel run each, against the same retrained model explained with two Kernel seeds.}
\label{tab:retrainseed}
\end{table}
\begin{table}[h]
\centering\small
\resizebox{\linewidth}{!}{\input{../figures/tab_refcheck.tex}}
\caption{German Credit reference check, per model: relative $\ell_2$ distance of the 20\,000-coalition reference (mean of 10 Kernel runs) to the exact vector by enumeration of $2^{20}$ coalitions (median and maximum over the instances checked), the reference's median relative standard error, and the number of instances whose top-3 and top-5 sets coincide.}
\label{tab:refcheck}
\end{table}
\begin{table}[h]
\centering\small
\resizebox{\linewidth}{!}{\input{../figures/tab_logit.tex}}
\caption{Explained output: probability of the positive class versus log-odds. Fraction of instances whose exact (reference) top-3 sets coincide, with 95\% CI, and the top-3 match probability to each output's own reference at the smallest and largest budget. For German logistic regression the log-odds game is linear, so Kernel SHAP is exact there. Only these five cells were run (ACS Employment gbm; German gbm, lr, mlp; Pima gbm).}
\label{tab:logit}
\end{table}

\section{Pre-specified choices}
\label{app:prereg}
The study plan is part of the delivery: \texttt{plan/study\_plan\_at\_launch.md} is the version in force when the main grid was launched, \texttt{plan/study\_plan.md} the final version, and \texttt{plan/CHANGES.md} lists every revision with its commit and time. The status of each analysis in this paper, read from those files and from the commit log of the analysis code, is in Table~\ref{tab:prereg}. ``Planned subset'' means the plan named the analysis but the paper's selection of cells or pairs was made later; ``post-hoc'' means the plan did not contain it. For the post-hoc items we cannot show whether they were written before or after the results were seen, so they should be read as exploratory.

\begin{table}[h]
\centering\small
\begin{tabular}{p{0.52\linewidth}p{0.4\linewidth}}
\toprule
Analysis & Status \\
\midrule
Measures (relative $\ell_2$, top-$k$ match and Jaccard, Kendall $\tau$, RBO, Nogueira, sign flips, margin); $k\in\{1,3,5\}$, $k=3$ for headlines & pre-specified \\
Population-level adequacy criterion, instance-cluster bootstrap, error-versus-budget slope fits & pre-specified \\
Headline cells: German Credit and ACS Employment at the library-default budget & pre-specified; ACS Employment added before any result \\
Comparison standards (retraining, background redraws, null) & pre-specified; run on a subset of cells \\
Default \texttt{l1\_reg} and log-odds output & pre-specified as sensitivity checks on a subset of cells; which cells chosen later \\
Per-instance diagnostic & pre-specified as optional, exploratory \\
Repeat versus enlarge & pre-specified as a comparison; pair selection (budget ratio within $[0.9,1.1]$) set after the budgets were fixed \\
50 instances (plan: 100), ACS Employment in the core, five retrain seeds, 50-tree forest, budgets not in ratio 4 & design changed at calibration, recorded after launch and before any result \\
Median and blow-up summaries, slope fits restricted to the Monte-Carlo-like range, skipping local slopes between budgets with equal rows, the gap-over-run-noise form of the margin analysis & post-hoc \\
Explainer disagreement, Covert--Lee variance comparison, Arrhythmia & planned, not run \\
\bottomrule
\end{tabular}
\caption{Status of the analyses relative to the study plan.}
\label{tab:prereg}
\end{table}
